\documentclass[11pt]{article}

\usepackage[T1]{fontenc}
\usepackage{lmodern}
\usepackage{microtype}
\usepackage[margin=1.08in]{geometry}
\usepackage{amsmath,amssymb,amsthm,mathtools}
\usepackage{booktabs}
\usepackage{enumitem}
\usepackage{xcolor}
\usepackage{listings}
\usepackage{hyperref}
\usepackage[nameinlink,noabbrev]{cleveref}
\crefname{theorem}{theorem}{theorems}
\Crefname{theorem}{Theorem}{Theorems}
\crefname{proposition}{proposition}{propositions}
\Crefname{proposition}{Proposition}{Propositions}
\crefname{lemma}{lemma}{lemmas}
\Crefname{lemma}{Lemma}{Lemmas}
\crefname{corollary}{corollary}{corollaries}
\Crefname{corollary}{Corollary}{Corollaries}
\crefname{remark}{remark}{remarks}
\Crefname{remark}{Remark}{Remarks}

\hypersetup{
  colorlinks=true,
  linkcolor=blue!55!black,
  citecolor=blue!55!black,
  urlcolor=blue!55!black,
  pdftitle={Audit of a Proposed Lorentzian Jensen--Bregman Metric Principle}
}

\newtheorem{theorem}{Theorem}[section]
\newtheorem{proposition}[theorem]{Proposition}
\newtheorem{lemma}[theorem]{Lemma}
\newtheorem{corollary}[theorem]{Corollary}
\newtheorem{remark}[theorem]{Remark}
\newtheorem{question}[theorem]{Question}
\theoremstyle{definition}
\newtheorem{definition}[theorem]{Definition}

\newcommand{\R}{\mathbb{R}}
\newcommand{\C}{\mathcal{C}}
\newcommand{\J}{\mathsf{J}}
\newcommand{\Vol}{\operatorname{Vol}}
\newcommand{\D}{\mathrm{D}}
\newcommand{\Q}{\mathcal{Q}}
\newcommand{\T}{\mathcal{T}}
\newcommand{\ip}[2]{\left\langle #1,#2\right\rangle}
\newcommand{\norm}[1]{\left\lVert #1\right\rVert}

\lstdefinestyle{pythonstyle}{
  language=Python,
  basicstyle=\ttfamily\small,
  columns=fullflexible,
  breaklines=true,
  frame=single,
  showstringspaces=false,
  keywordstyle=\color{blue!60!black},
  commentstyle=\color{green!40!black},
  stringstyle=\color{red!55!black}
}

\title{Counterexamples to a Proposed Lorentzian\\Jensen--Bregman Metric Principle}
\author{}
\date{\today}

\begin{document}
\maketitle

\begin{abstract}
We audit the proposed principle that, for every positive homogeneous Lorentzian polynomial
$G$, the midpoint Jensen gap generated by $F=-\log G$ is the square of a metric.  The
proposed proof proceeds through a claimed Gram-determinant formula for the derivative of
a Hessian Rayleigh quotient.  We show that this derivative formula is impossible on
homogeneity and parity grounds and fails even for the volume polynomial of elementary
planar rectangles.  We then give an explicit strictly Lorentzian bivariate cubic and three
rational points for which the square root of the midpoint Jensen gap violates the triangle
inequality; the separation is certified by exact rational enclosures for logarithms and square
roots.  Finally, Shephard's realization theorem for the mixed volumes of two convex
bodies transfers the same counterexample to the proposed ``Jensen--Shannon volume''
distance on convex bodies in $\R^3$.  We record the exact derivative identity, isolate the
valid parts of the earlier argument, and formulate a corrected research agenda.
\end{abstract}

\noindent\textbf{Keywords.} Jensen divergence; Lorentzian polynomial; mixed volume; Steiner polynomial; triangle inequality; counterexample.

\medskip
\noindent\textbf{2020 Mathematics Subject Classification.} 52A39, 26B25, 15A45, 05E40.

\section{The proposed principle and the audit verdict}

Let $\C\subset\R^m$ be an open convex cone, let $G:\C\to(0,\infty)$, and put
\[
 F=-\log G,
 \qquad
 \J_G(x,y)
 :=\frac{F(x)+F(y)}2-F\!\left(\frac{x+y}{2}\right)
 =\log G\!\left(\frac{x+y}{2}\right)
   -\frac{\log G(x)+\log G(y)}2.
\]
The proposed master theorem asserted that if $G$ is homogeneous and Lorentzian, then
$d_G:=\sqrt{\J_G}$ is a semimetric.  A geometric specialization asserted that
\[
 d_{\Vol}(K,L)
 :=\left[
 \log\Vol\!\left(\frac{K+L}{2}\right)
 -\frac12\log\Vol(K)-\frac12\log\Vol(L)
 \right]^{1/2}
\]
is a semimetric on convex bodies.

\medskip
\noindent
\textbf{Audit verdict.}
\begin{enumerate}[label=\textup{(\roman*)},leftmargin=2.1em]
\item The midpoint integral identity and the homogeneous primitive decomposition used in the
      proposed proof are correct.
\item The claimed Gram-determinant identity for the derivative of a Hessian Rayleigh quotient
      is false.  Its two sides have incompatible homogeneity in the line direction.
\item The implication ``Lorentzian $G$ $\Rightarrow$ $\sqrt{\J_G}$ is a metric'' is false,
      already for a strictly Lorentzian bivariate cubic of degree three.
\item The convex-body specialization is also false: by Shephard's theorem the same cubic is
      the volume polynomial of two three-dimensional convex bodies.
\item Consequently, the earlier positive paper cannot be made rigorous by filling in the
      omitted Gram-determinant calculation; the omitted calculation conceals a genuine false
      identity rather than a merely technical gap.
\end{enumerate}

\section{Correct preliminary identities}

\subsection{The midpoint integral identity}

\begin{lemma}[Midpoint Green-kernel identity]\label{lem:midpoint}
Let $F\in C^2(\C)$ be convex, and let $x,y\in\C$.  Set
$m=(x+y)/2$ and $h=y-x$.  Then
\begin{equation}\label{eq:midpoint}
 \frac{F(x)+F(y)}2-F(m)
 =\frac18\int_{-1}^{1}(1-|t|)
   \,\D^2F\!\left(m+\frac{t}{2}h\right)[h,h] \,dt.
\end{equation}
\end{lemma}

\begin{proof}
Define $\varphi(s)=F(m+sh)$ for $s\in[-1/2,1/2]$.  For $a\ge0$,
Taylor's formula with integral remainder gives
\[
 \varphi(a)=\varphi(0)+a\varphi'(0)+\int_0^a(a-u)\varphi''(u)\,du
\]
and, after changing orientation in the integral,
\[
 \varphi(-a)=\varphi(0)-a\varphi'(0)+\int_0^a(a-u)\varphi''(-u)\,du.
\]
Taking $a=1/2$, adding, and dividing by two yields
\[
 \frac{\varphi(1/2)+\varphi(-1/2)}2-\varphi(0)
 =\frac12\int_0^{1/2}\left(\frac12-u\right)
   \bigl(\varphi''(u)+\varphi''(-u)\bigr)\,du.
\]
Using $\varphi''(u)=\D^2F(m+uh)[h,h]$ and the substitutions $t=2u$ and
$t=-2u$ gives \eqref{eq:midpoint}.
\end{proof}

\subsection{The exact directional derivative}

For a fixed vector $u\in\R^m$, define the Hessian Rayleigh quotient
\[
 q_u(r):=\D^2F(r)[u,u],\qquad F=-\log G.
\]
The following identity is the correct starting point for any monotonicity argument.

\begin{proposition}[Exact third-order derivative identity]\label{prop:direct-derivative}
Let $G\in C^3(\C)$ be positive.  For every $r\in\C$ and $u,\delta\in\R^m$,
\begin{align}
 \D_\delta q_u(r)
 &= -\D^3(\log G)(r)[u,u,\delta] \label{eq:third-log}\\
 &= \frac{2G_uG_{u\delta}}{G^2}
    -\frac{2G_u^2G_\delta}{G^3}
    -\frac{G_{uu\delta}}{G}
    +\frac{G_{uu}G_\delta}{G^2},               \label{eq:direct-expanded}
\end{align}
where every derivative on the right is evaluated at $r$ and
$G_u=\D_uG$, $G_{u\delta}=\D^2G[u,\delta]$, and so forth.
\end{proposition}

\begin{proof}
Since $F=-\log G$ and directional derivatives commute,
\[
 \D_\delta q_u
 =\D_\delta\D_u^2F
 =-\D_\delta\D_u^2\log G
 =-\D_u^2\D_\delta\log G,
\]
which is \eqref{eq:third-log}.  Moreover,
\[
 q_u=\D_u^2(-\log G)=\frac{G_u^2}{G^2}-\frac{G_{uu}}G.
\]
Differentiating the first term gives
\[
 \D_\delta\!\left(\frac{G_u^2}{G^2}\right)
 =\frac{2G_uG_{u\delta}}{G^2}
  -\frac{2G_u^2G_\delta}{G^3},
\]
while differentiating the second gives
\[
 \D_\delta\!\left(-\frac{G_{uu}}G\right)
 =-\frac{G_{uu\delta}}G+\frac{G_{uu}G_\delta}{G^2}.
\]
Adding proves \eqref{eq:direct-expanded}.
\end{proof}

\begin{corollary}[Oddness and linear homogeneity]\label{cor:oddness}
For fixed $r$ and $u$, the map $\delta\mapsto\D_\delta q_u(r)$ is linear.  In particular,
\[
 \D_{c\delta}q_u(r)=c\,\D_\delta q_u(r),
 \qquad
 \D_{-\delta}q_u(r)=-\D_\delta q_u(r).
\]
\end{corollary}

\subsection{Homogeneous primitive coordinates}

Assume now that $G$ is homogeneous of degree $d\ge2$.  Define
\[
 \Q_r(v,w):=\frac{1}{d(d-1)}\D^2G(r)[v,w],
 \qquad
 \T_r(v,w,z):=\frac{1}{d(d-1)}\D^3G(r)[v,w,z].
\]
Euler's identities imply
\begin{equation}\label{eq:euler-normalized}
 \Q_r(r,r)=G(r),
 \qquad
 \Q_r(v,r)=\frac1dG_v(r),
 \qquad
 \T_r(v,w,r)=(d-2)\Q_r(v,w).
\end{equation}
For $u,\delta\in\R^m$, set
\[
 \alpha=\frac{\Q_r(u,r)}{G(r)},\qquad
 \beta =\frac{\Q_r(\delta,r)}{G(r)},\qquad
 U=u-\alpha r,\qquad V=\delta-\beta r.
\]
Then $\Q_r(U,r)=\Q_r(V,r)=0$.

\begin{lemma}[Primitive decomposition]\label{lem:primitive-decomp}
With the notation above,
\begin{equation}\label{eq:primitive-decomp}
 q_u(r)
 =\frac1d\left(\frac{G_u(r)}{G(r)}\right)^2
  -\frac{d(d-1)}{G(r)}\Q_r(U,U).
\end{equation}
\end{lemma}

\begin{proof}
By \eqref{eq:euler-normalized}, $\alpha=G_u/(dG)$.  Expanding gives
\[
 \Q_r(U,U)=\Q_r(u,u)-2\alpha\Q_r(u,r)+\alpha^2\Q_r(r,r)
 =\Q_r(u,u)-\frac{G_u^2}{d^2G}.
\]
Thus
\[
 \D^2G[u,u]=d(d-1)\Q_r(u,u)
 =d(d-1)\Q_r(U,U)+\frac{d-1}{d}\frac{G_u^2}{G}.
\]
Substitution into
$q_u=G_u^2/G^2-\D^2G[u,u]/G$ yields \eqref{eq:primitive-decomp}.
\end{proof}

The next proposition records the exact primitive-coordinate derivative.  It is included to make
fully explicit the term that the proposed proof incorrectly replaced by a Gram determinant.

\begin{proposition}[Exact derivative in primitive coordinates]\label{prop:primitive-derivative}
With $S=G(r)$ and the notation above,
\begin{align}
 \D_\delta q_u(r)
 ={}&-\frac{d(d-1)}S\T_r(U,U,V)
     +\frac{4d(d-1)\alpha}{S}\Q_r(U,V) \label{eq:primitive-derivative}\\
    &+\frac{2d(d-1)\beta}{S}\Q_r(U,U)
     -2d\alpha^2\beta.\nonumber
\end{align}
\end{proposition}

\begin{proof}
Write, at the fixed point $r$,
\[
 a=G_u=d\alpha S,
 \quad b=G_\delta=d\beta S,
 \quad c=G_{uu}=d(d-1)\Q_r(u,u),
\]
\[
 e=G_{u\delta}=d(d-1)\Q_r(u,\delta),
 \quad t=G_{uu\delta}=d(d-1)\T_r(u,u,\delta).
\]
By \eqref{eq:direct-expanded},
\begin{equation}\label{eq:four-terms}
 \D_\delta q_u=\frac{2ae}{S^2}-\frac{2a^2b}{S^3}-\frac{t}{S}+\frac{cb}{S^2}.
\end{equation}
Because $u=U+\alpha r$ and $\delta=V+\beta r$, primitiveness gives
\begin{align}
 \Q_r(u,\delta)&=\Q_r(U,V)+\alpha\beta S,\label{eq:Qud}\\
 \Q_r(u,u)&=\Q_r(U,U)+\alpha^2S.\label{eq:Quu}
\end{align}
By multilinearity, symmetry, and \eqref{eq:euler-normalized},
\begin{align}
 \T_r(u,u,\delta)
 &=\T_r(U+\alpha r,U+\alpha r,V+\beta r)\nonumber\\
 &=\T_r(U,U,V)
   +\beta\T_r(U,U,r)
   +2\alpha\T_r(U,r,V)
   +2\alpha\beta\T_r(U,r,r)\nonumber\\
 &\qquad
   +\alpha^2\T_r(r,r,V)
   +\alpha^2\beta\T_r(r,r,r)\nonumber\\
 &=\T_r(U,U,V)
   +\beta(d-2)\Q_r(U,U)
   +2\alpha(d-2)\Q_r(U,V)
   +\alpha^2\beta(d-2)S.                 \label{eq:Texpand}
\end{align}
Here $\T_r(U,r,r)=(d-2)\Q_r(U,r)=0$ and similarly
$\T_r(r,r,V)=0$.

We now expand each term in \eqref{eq:four-terms}.  Using \eqref{eq:Qud},
\begin{align*}
 \frac{2ae}{S^2}
 &=\frac{2(d\alpha S)d(d-1)}{S^2}
   \bigl(\Q_r(U,V)+\alpha\beta S\bigr)\\
 &=\frac{2d^2(d-1)\alpha}{S}\Q_r(U,V)
   +2d^2(d-1)\alpha^2\beta.
\end{align*}
Also,
\[
 -\frac{2a^2b}{S^3}=-2d^3\alpha^2\beta.
\]
Using \eqref{eq:Texpand},
\begin{align*}
 -\frac{t}{S}
 ={}&-\frac{d(d-1)}S\T_r(U,U,V)
     -\frac{d(d-1)(d-2)\beta}{S}\Q_r(U,U)\\
    &-\frac{2d(d-1)(d-2)\alpha}{S}\Q_r(U,V)
     -d(d-1)(d-2)\alpha^2\beta.
\end{align*}
Finally, by \eqref{eq:Quu},
\[
 \frac{cb}{S^2}
 =\frac{d^2(d-1)\beta}{S}\Q_r(U,U)
  +d^2(d-1)\alpha^2\beta.
\]
Collecting the coefficients of $\Q_r(U,V)$ gives
\[
 \frac{2d(d-1)\alpha}{S}\bigl(d-(d-2)\bigr)
 =\frac{4d(d-1)\alpha}{S}.
\]
Collecting the coefficients of $\Q_r(U,U)$ gives
\[
 \frac{d(d-1)\beta}{S}\bigl(d-(d-2)\bigr)
 =\frac{2d(d-1)\beta}{S}.
\]
The coefficient of $\alpha^2\beta$ equals
\begin{align*}
 &2d^2(d-1)-2d^3-d(d-1)(d-2)+d^2(d-1)\\
 &\qquad
 =d\bigl(3d(d-1)-2d^2-(d-1)(d-2)\bigr)=-2d.
\end{align*}
This proves \eqref{eq:primitive-derivative}.
\end{proof}

\section{Why the claimed Gram-determinant formula cannot hold}

The proposed formula had the schematic form
\begin{equation}\label{eq:false-schematic}
 \D_\delta q_u(r)
 =C(r)\det
 \begin{pmatrix}
  \langle U,U\rangle_r&\langle U,V\rangle_r\\
  \langle V,U\rangle_r&\langle V,V\rangle_r
 \end{pmatrix},
\end{equation}
where $V$ is the primitive projection of $\delta$ and the inner product is derived from
$-\Q_r$.

\begin{proposition}[Homogeneity obstruction]\label{prop:homogeneity-obstruction}
No identity of the form \eqref{eq:false-schematic}, with $C(r)$ independent of $\delta$,
can hold for all line directions $\delta$, except in the trivial case in which both sides vanish.
\end{proposition}

\begin{proof}
By Corollary~\ref{cor:oddness}, the left side is homogeneous of degree one and odd in $\delta$.
The primitive projection $V$ is linear in $\delta$, so the Gram determinant on the right is
homogeneous of degree two and even in $\delta$.  Replacing $\delta$ by $c\delta$ therefore
multiplies the two sides by $c$ and $c^2$, respectively; replacing $\delta$ by $-\delta$
changes only the sign of the left side.  Thus the identity is impossible unless both sides
are identically zero.
\end{proof}

\begin{remark}[Orientation obstruction for the third-order term]\label{rem:T-sign}
The same parity observation rules out the proposed shortcut
$\T_r(U,U,V)\le0$ for every line direction.  For fixed $r$ and $u$, replacing
$\delta$ by $-\delta$ replaces its primitive projection $V$ by $-V$ and hence changes the
sign of $\T_r(U,U,V)$.  A one-sided sign valid for both orientations would therefore force
$\T_r(U,U,V)=0$ identically.  Likewise, an assertion
$\D_\delta q_u(r)\ge0$ for every $\delta$ would, after replacing $\delta$ by $-\delta$,
force every directional derivative to vanish.  Any valid monotonicity statement must therefore
select an oriented cone of directions and cannot follow from an even Gram determinant.
\end{remark}

\subsection{An explicit mixed-volume contradiction}

The failure is visible even in dimension two, with full-dimensional axis-aligned rectangles.
Let
\[
 X_a=X_c=[0,1]\times[0,1],
 \qquad
 X_b=[0,2]\times[0,3].
\]
Fix $\theta=1/2$ and, for $0<s<1/2$, put
\[
 r(s)=(1-\theta)e_a+s e_b+(\theta-s)e_c.
\]
Since $X_a=X_c$,
\[
 M(s)=\sum_i r_i(s)X_i=(1-s)X_a+sX_b
      =[0,1+s]\times[0,1+2s],
\]
and hence
\[
 G(r(s))=\Vol(M(s))=(1+s)(1+2s).
\]
For $u=e_b-e_a$ and $F=-\log G$, the two side-length functionals have directional
increments $1$ and $2$, respectively.  Therefore
\begin{equation}\label{eq:rectangle-q}
 q_u(r(s))=\D^2F(r(s))[u,u]
 =\frac{1}{(1+s)^2}+\frac{4}{(1+2s)^2},
\end{equation}
and
\begin{equation}\label{eq:rectangle-qprime}
 \frac{d}{ds}q_u(r(s))
 =-\frac{2}{(1+s)^3}-\frac{16}{(1+2s)^3}<0.
\end{equation}
The proposed mixed-volume Gram formula carried a prefactor $n(n-1)(n-2)$ and would
therefore give zero when $n=2$, contradicting \eqref{eq:rectangle-qprime}.  Thus the
failure is not specific to arbitrary Lorentzian polynomials; the stated mixed-volume
calculation itself was false.

\section{A strictly Lorentzian counterexample to metricity}

We now disprove the proposed master theorem itself.

\begin{theorem}[Lorentzian cubic with a triangle violation]\label{thm:lorentzian-counterexample}
Let
\begin{equation}\label{eq:cubic}
 G(x,y)=y^3+\frac{11}{5}xy^2+\frac32x^2y+\frac1{10}x^3,
 \qquad (x,y)\in\R_{>0}^2,
\end{equation}
and let $F=-\log G$.  Then $G$ is strictly Lorentzian, but
$d_F=\sqrt{\J_G}$ does not satisfy the triangle inequality on $\R_{>0}^2$.
\end{theorem}

\begin{proof}
Write $G(x,y)=\sum_{k=0}^3 a_kx^ky^{3-k}$, so
\[
 (a_0,a_1,a_2,a_3)=\left(1,\frac{11}{5},\frac32,\frac1{10}\right).
\]
The bivariate Lorentzian criterion \cite[Example~2.3]{BrandenHuh} states that strict Lorentzianity is equivalent to strict
log-concavity of the binomially normalized coefficients $a_k/\binom3k$.  Here
\[
 \left(\frac{a_0}{\binom30},\frac{a_1}{\binom31},
       \frac{a_2}{\binom32},\frac{a_3}{\binom33}\right)
 =\left(1,\frac{11}{15},\frac12,\frac1{10}\right),
\]
and
\[
 \left(\frac{11}{15}\right)^2=\frac{121}{225}>\frac12,
 \qquad
 \left(\frac12\right)^2=\frac14>\frac{11}{150}.
\]
Thus $G$ is strictly Lorentzian.  In particular, $G$ is log-concave on the positive
orthant \cite{BrandenHuh}, so $F=-\log G$ is convex and all the Jensen gaps below are
nonnegative.

Consider the rational points
\begin{equation}\label{eq:points}
 p=(35,2),
 \qquad
 q=\left(\frac{11}{2},\frac{23}{2}\right),
 \qquad
 r=\left(\frac1{500},15\right).
\end{equation}
A direct exact evaluation gives
\begin{center}
\begin{tabular}{@{}c@{\qquad}c@{}}
\toprule
point $z$ & $G(z)$\\
\midrule
$p$ & $16557/2$\\
$q$ & $73191/20$\\
$r$ & $4219987612501/1250000000$\\
$(p+q)/2$ & $2342277/320$\\
$(q+r)/2$ & $35412384644501/10000000000$\\
$(p+r)/2$ & $78370720930001/10000000000$\\
\bottomrule
\end{tabular}
\end{center}
Consequently, if
\[
 R_{ab}:=\frac{G((a+b)/2)^2}{G(a)G(b)},
 \qquad
 \J_G(a,b)=\frac12\log R_{ab},
\]
then
\begin{align*}
 R_{pq}&=\frac{609584616081}{344696430080},\\
 R_{qr}&=\frac{1254036986210090216149539001}
               {1235460453386242764000000000},\\
 R_{pr}&=\frac{6141969899088096806341860001}
               {2794813396007162280000000000}.
\end{align*}
The exact rational enclosure described in \cref{app:certificate} gives
\begin{align*}
 \J_G(p,q)&=0.28505682554857745905016398352381466439\ldots,\\
 \J_G(q,r)&=0.00746209950879334567961573595521132391\ldots,\\
 \J_G(p,r)&=0.39369009014987658353757866710396789908\ldots.
\end{align*}
More importantly, using only rational arithmetic and integer square roots, the certificate proves
\begin{equation}\label{eq:exact-violation-lower-bound}
 d_F(p,r)-d_F(p,q)-d_F(q,r)
 >\frac{7156704176613847159228106087505875849130467}{10^{45}}>0.
\end{equation}
Thus the triangle inequality fails.
\end{proof}

\section{The convex-body distance is not a metric}

For convex bodies $K,L\subset\R^3$, the relative Steiner polynomial has the form
\begin{equation}\label{eq:steiner}
 \Vol(xK+yL)
 =\sum_{i=0}^3\binom3iW_i(K,L)x^{3-i}y^i.
\end{equation}
Shephard proved that, for two convex bodies, the Alexandrov--Fenchel inequalities are a
full set \cite[Theorem~4]{Shephard}; see also \cite[Lemma~2.1]{HenkHernandezSaorin}: every positive sequence satisfying the adjacent log-concavity inequalities occurs as
the mixed-volume sequence of some pair of convex bodies.

\begin{corollary}[Counterexample on convex bodies]\label{cor:convex-body-counterexample}
There exist full-dimensional convex bodies $K,L\subset\R^3$ and three Minkowski
combinations $A,B,C$ of $K,L$ for which
\[
 d_{\Vol}(A,C)>d_{\Vol}(A,B)+d_{\Vol}(B,C).
\]
In particular, the proposed Jensen--volume square root is not a semimetric on convex bodies.
\end{corollary}

\begin{proof}
Set
\[
 (W_0,W_1,W_2,W_3)
 =\left(\frac1{10},\frac12,\frac{11}{15},1\right).
\]
The complete set of Shephard inequalities in dimension three is satisfied; explicitly,
\[
 W_1^2=\frac14>\frac{11}{150}=W_0W_2,
 \qquad
 W_1W_2=\frac{11}{30}>\frac1{10}=W_0W_3,
 \qquad
 W_2^2=\frac{121}{225}>\frac12=W_1W_3.
\]
By Shephard's realization theorem \cite[Theorem~4]{Shephard}, there exist convex bodies $K,L\subset\R^3$ whose
mixed-volume sequence is the displayed tuple.  Their polynomial \eqref{eq:steiner} is
therefore exactly the cubic \eqref{eq:cubic}:
\[
 \Vol(xK+yL)=\frac1{10}x^3+\frac32x^2y+\frac{11}{5}xy^2+y^3=G(x,y).
\]
The positive endpoint coefficients imply $\Vol(K),\Vol(L)>0$, so both bodies are
full-dimensional.

Let $p,q,r$ be as in \eqref{eq:points} and define
\[
 A=p_1K+p_2L,
 \qquad B=q_1K+q_2L,
 \qquad C=r_1K+r_2L.
\]
Minkowski linearity gives
\[
 \frac{A+B}{2}=\frac{p_1+q_1}{2}K+\frac{p_2+q_2}{2}L,
\]
and similarly for the other pairs.  Hence
\[
 d_{\Vol}(A,B)^2=\J_G(p,q),
 \quad d_{\Vol}(B,C)^2=\J_G(q,r),
 \quad d_{\Vol}(A,C)^2=\J_G(p,r).
\]
The strict triangle violation in \cref{thm:lorentzian-counterexample} transfers verbatim.
\end{proof}


\section{Status of the matrix three-lines component}

The abstract interpolation statement and its proposed geometric application must be separated.
The following points survive the audit.

\begin{enumerate}[label=\textup{(\roman*)},leftmargin=2.1em]
\item The original Gram statement with an analytic family $T_z$ but endpoint Grams formed by
      manually replacing $T_0,T_1$ by the identity is false: the endpoint data need not be the
      boundary values of the analytic family.
\item If one instead defines the endpoint maps from the actual boundary values and assumes
      contractivity in the corresponding weighted Hilbert norms, ordinary complex interpolation
      of Hilbert couples does give a valid Gram inequality.  The interpolated domain form is the
      Kubo--Ando geometric mean of the two endpoint Gram forms (with the usual regularization in
      the singular case).
\item This corrected abstract lemma does not prove the desired Fisher-tensor inequality.  The
      analytic family built from left and right matrix powers acts on a common seed, whereas the
      three Fisher Grams in the geometric argument are built from different normalized seeds.
      There is no general identification of those endpoint Grams.
\item A scalar trace-product estimate is not the dual formulation of the Loewner comparison; the
      valid dual test must pair the test matrix with the geometric mean itself.
\end{enumerate}

Thus the matrix three-lines lemma is not intrinsically false after its hypotheses and endpoint maps
are corrected, but it does not bridge the noncommutative normalization mismatch in the proposed
proof.

\section{Additional defects in the earlier proof chain}

For completeness, we record three further logical gaps that are independent of the explicit
counterexample above.

\subsection{Changing probability measures under H\"older}

Suppose a pointwise estimate has the form
\[
 \Psi_{ac}(\zeta)\le \Psi_{ab}(\zeta)^{1-\theta}\Psi_{bc}(\zeta)^\theta.
\]
Integrating against $\mu_{ac}$ and applying H\"older yields only
\[
 \int\Psi_{ac}\,d\mu_{ac}
 \le
 \left(\int\Psi_{ab}\,d\mu_{ac}\right)^{1-\theta}
 \left(\int\Psi_{bc}\,d\mu_{ac}\right)^\theta.
\]
It does not produce integrals against $\mu_{ab}$ and $\mu_{bc}$.  A change of measure would
require explicit Radon--Nikodym factors and new estimates.  Thus the earlier ``integrated scalar
GM'' theorem did not follow from its stated proof.

\subsection{The covariance subtraction has the wrong monotonicity direction}

Writing $H_\alpha=\mathsf F_\alpha-\operatorname{Cov}_\alpha$ gives
$0\preceq H_\alpha\preceq\mathsf F_\alpha$.  Even if one proved
\[
 H_{ac}\preceq \mathsf F_{ac}
 \preceq \mathsf F_{ab}\#_\theta\mathsf F_{bc},
\]
one could not replace the endpoint Fisher tensors by $H_{ab}$ and $H_{bc}$.  Monotonicity of
the Kubo--Ando mean gives the opposite comparison,
\[
 H_{ab}\#_\theta H_{bc}
 \preceq \mathsf F_{ab}\#_\theta\mathsf F_{bc}.
\]
Thus the desired smaller upper bound does not follow.

\subsection{The trace-product test is weaker than Loewner order}

For positive semidefinite $A,B,C$,
\[
 \operatorname{tr}\bigl(C(A\#_\theta B)\bigr)
 \le
 \operatorname{tr}(CA)^{1-\theta}\operatorname{tr}(CB)^\theta.
\]
Therefore an estimate of the form
\[
 \operatorname{tr}(CX)
 \le \operatorname{tr}(CA)^{1-\theta}\operatorname{tr}(CB)^\theta
 \qquad(C\succeq0)
\]
does not imply $X\preceq A\#_\theta B$: its right side is generally larger than the dual
pairing with $A\#_\theta B$.  The valid dual characterization is instead
$X\preceq Y$ if and only if $\operatorname{tr}(CX)\le\operatorname{tr}(CY)$ for every
$C\succeq0$.

\section{What survives and a corrected research direction}

The counterexamples do not affect the following established or elementary facts.

\begin{enumerate}[label=\textup{(\alph*)},leftmargin=2em]
\item In one dimension, Chen--Chen--Rao \cite{ChenChenRao} characterize metricity by
      convexity of $\log F''$ (under the regularity and positivity hypotheses of their theorem).
\item The same result extends to separable sums and, more generally, to nonnegative
      superpositions of one-dimensional generators.
\item The Jensen--Bregman log-determinant divergence is the S-divergence (also called the
      Stein divergence); its square root is a known metric \cite{SraSdiv}.  That fact does not
      follow from Lorentzianity alone.
\item Any genuine higher-dimensional characterization must control arbitrary triangles;
      linewise log-convexity is necessary but, by itself, is not presently known to be sufficient.
\end{enumerate}

The following elementary closure principle is a useful positive substitute for the false
Lorentzian theorem.

\begin{proposition}[Nonnegative superposition principle]\label{prop:superposition}
Let $(\Omega,\mu)$ be a measure space.  For each $\omega$, let $I_\omega\subset\R$ be an
interval, let $f_\omega:I_\omega\to\R$ be convex, and let
$\ell_\omega:\C\to I_\omega$ be affine.  Assume that
$d_\omega(s,t):=\sqrt{\J_{f_\omega}(s,t)}$ is a semimetric on $I_\omega$ and that
\[
 F(x):=\int_\Omega f_\omega(\ell_\omega(x))\,d\mu(\omega)
\]
is finite.  Then $\sqrt{\J_F}$ is a semimetric on $\C$.
\end{proposition}

\begin{proof}
Affineness of $\ell_\omega$ and linearity of the Jensen gap give
\[
 \J_F(x,y)=\int_\Omega
 \J_{f_\omega}(\ell_\omega(x),\ell_\omega(y))\,d\mu(\omega)
 =\int_\Omega d_\omega(\ell_\omega(x),\ell_\omega(y))^2\,d\mu(\omega).
\]
For arbitrary $x,y,z$, the pointwise triangle inequality for $d_\omega$, followed by
Minkowski's inequality in $L^2(\mu)$, gives
\begin{align*}
 \sqrt{\J_F(x,z)}
 &\le\left(\int_\Omega
   \bigl[d_\omega(\ell_\omega(x),\ell_\omega(y))
        +d_\omega(\ell_\omega(y),\ell_\omega(z))\bigr]^2d\mu\right)^{1/2}\\
 &\le \sqrt{\J_F(x,y)}+\sqrt{\J_F(y,z)}.
\end{align*}
\end{proof}

A necessary higher-dimensional condition follows immediately by restriction to lines.

\begin{proposition}[Linewise necessary condition]\label{prop:linewise-necessary}
Assume $F\in C^4(\C)$ is strictly convex and $\sqrt{\J_F}$ is a metric on $\C$.
For every $x\in\C$ and every nonzero $h$ such that $x+th\in\C$ on an interval $I$,
\[
 t\longmapsto \log\bigl(\D^2F(x+th)[h,h]\bigr)
\]
is convex on $I$.
\end{proposition}

\begin{proof}
The restriction $f(t)=F(x+th)$ has midpoint Jensen distance equal to the restriction of
$\sqrt{\J_F}$ to the affine line.  It is therefore a metric on $I$.  The one-dimensional
necessity theorem of Chen--Chen--Rao \cite{ChenChenRao} gives convexity of $\log f''$, and
$f''(t)=\D^2F(x+th)[h,h]$.
\end{proof}

\begin{question}
Find a genuinely nonseparable higher-dimensional condition that is both stable under useful
operations and sufficient for metricity of $\sqrt{\J_F}$.  The counterexample in
\cref{thm:lorentzian-counterexample} shows that complete log-concavity/Lorentzianity of
$G=e^{-F}$ is not such a condition.
\end{question}

\appendix
\section{Exact rational certificate for the triangle violation}\label{app:certificate}

The numerical separation in \eqref{eq:exact-violation-lower-bound} can be certified without
floating-point assumptions.  We use the following elementary enclosure.

\begin{lemma}[Rational enclosure for the logarithm]\label{lem:log-enclosure}
Let $x>1$ be rational, set $z=(x-1)/(x+1)\in(0,1)$, and let $N\ge1$.  Then
\[
 L_N(x):=2\sum_{k=0}^{N-1}\frac{z^{2k+1}}{2k+1}
 <\log x
 <L_N(x)+\frac{2z^{2N+1}}{(2N+1)(1-z^2)}.
\]
Every quantity in this enclosure is rational.
\end{lemma}

\begin{proof}
Since $x=(1+z)/(1-z)$,
\[
 \log x=2\operatorname{artanh}z
 =2\sum_{k=0}^{\infty}\frac{z^{2k+1}}{2k+1}.
\]
All terms are positive.  For $k\ge N$, $(2k+1)^{-1}\le(2N+1)^{-1}$, and hence the tail is
bounded above by
\[
 \frac{2}{2N+1}\sum_{k=N}^{\infty}z^{2k+1}
 =\frac{2z^{2N+1}}{(2N+1)(1-z^2)}.
\]
\end{proof}

For a positive rational $q$ and a decimal scale $S=10^P$, the integer
\[
 m=\left\lfloor\sqrt{qS^2}\right\rfloor
\]
gives the exact enclosure $m/S\le\sqrt q<(m+1)/S$.  The script below combines this
observation with Lemma~\ref{lem:log-enclosure}, taking $N=100$ and $P=45$.  It uses only
\texttt{Fraction} and the exact integer square-root routine \texttt{isqrt}; its final assertion is
therefore an exact rational check.

\begin{lstlisting}[style=pythonstyle]
from fractions import Fraction
from math import isqrt


def G(point):
    x, y = point
    return (y**3 + Fraction(11, 5)*x*y**2
            + Fraction(3, 2)*x**2*y + Fraction(1, 10)*x**3)


def midpoint(a, b):
    return tuple((a[i] + b[i])/2 for i in range(2))


def ratio(a, b):
    return G(midpoint(a, b))**2 / (G(a)*G(b))


def log_interval(x, n_terms=100):
    if x <= 0:
        raise ValueError("x must be positive")
    if x == 1:
        return Fraction(0), Fraction(0)
    if x < 1:
        lower, upper = log_interval(1/x, n_terms)
        return -upper, -lower

    z = (x - 1)/(x + 1)
    z2 = z*z
    power = z
    partial = Fraction(0)
    for k in range(n_terms):
        partial += 2*power/Fraction(2*k + 1)
        power *= z2

    remainder = 2*power/(Fraction(2*n_terms + 1)*(1 - z2))
    return partial, partial + remainder


def sqrt_decimal_interval(x, digits=45):
    scale = 10**digits
    floor_argument = (x.numerator*scale*scale)//x.denominator
    m = isqrt(floor_argument)
    return Fraction(m, scale), Fraction(m + 1, scale)


p = (Fraction(35), Fraction(2))
q = (Fraction(11, 2), Fraction(23, 2))
r = (Fraction(1, 500), Fraction(15))

Lpq, Upq = log_interval(ratio(p, q))
Lqr, Uqr = log_interval(ratio(q, r))
Lpr, Upr = log_interval(ratio(p, r))

Jpq = (Lpq/2, Upq/2)
Jqr = (Lqr/2, Uqr/2)
Jpr = (Lpr/2, Upr/2)

sqrt_pr_lower, _ = sqrt_decimal_interval(Jpr[0])
_, sqrt_pq_upper = sqrt_decimal_interval(Jpq[1])
_, sqrt_qr_upper = sqrt_decimal_interval(Jqr[1])

violation_lower = sqrt_pr_lower - sqrt_pq_upper - sqrt_qr_upper
print(violation_lower)
assert violation_lower > 0
\end{lstlisting}

The printed rational number is
\[
 \frac{7156704176613847159228106087505875849130467}{10^{45}},
\]
which proves \eqref{eq:exact-violation-lower-bound}.

\begin{thebibliography}{99}

\bibitem{BrandenHuh}
P.~Br\"and\'en and J.~Huh,
\newblock Lorentzian polynomials,
\newblock \emph{Ann. of Math. (2)} \textbf{192} (2020), 821--891.

\bibitem{ChenChenRao}
P.~Chen, Y.~Chen, and M.~Rao,
\newblock Metrics defined by Bregman divergences,
\newblock \emph{Commun. Math. Sci.} \textbf{6} (2008), 915--926.

\bibitem{HenkHernandezSaorin}
M.~Henk, M.~A.~Hern\'andez Cifre, and E.~Saor\'in,
\newblock Steiner polynomials via ultra-logconcave sequences,
\newblock \emph{Commun. Contemp. Math.} \textbf{14} (2012), no.~6, 1250040.

\bibitem{Shephard}
G.~C.~Shephard,
\newblock Inequalities between mixed volumes of convex sets,
\newblock \emph{Mathematika} \textbf{7} (1960), 125--138.

\bibitem{SraSdiv}
S.~Sra,
\newblock Positive definite matrices and the S-divergence,
\newblock \emph{Proc. Amer. Math. Soc.} \textbf{144} (2016), 2787--2797.

\end{thebibliography}

\end{document}
